Let $c_n : \remainders{n} \to \remainders{n}$, $c_n(x) := -x - 1$.

\begin{prop}\label{prop.twosComplement}
 For every $x \in \remainders{n}$,
 $\digitsOf{n}(c_n(x)) = \neg\, \digitsOf{n}(x)$:
 complementing the $n+1$ digits of $x$ computes $-x-1$ modulo $2^{n+1}$.
\end{prop}
\begin{proof}
 Write $\digitsOf{n}(x) = (f_0, \dots, f_n)$, so that $x = \sum_{k=0}^{n} f_k 2^k$.
 The complemented digits have the value
 $\sum_{k=0}^{n} (1 - f_k) 2^k = (2^{n+1} - 1) - x$,
 which is $-x - 1$ modulo $2^{n+1}$.
\end{proof}

Equivalently, $\neg x + 1 = -x$ in $\remainders{n}$,
which is all that the rule ``invert the digits and add one'' says.
Listing \ref{listing.twosComplement} checks it on every remainder of ten moduli.

\begin{lstlisting}[caption={Complementing the digits computes $-x-1$.}, label=listing.twosComplement]
for n in range(10):
    m = 2 ** (n + 1)
    for x in range(m):
        assert (-x - 1) % m == (~x) % m == (m - 1) - x
\end{lstlisting}
